% =====================================================================
% P7 — Multi-country test of the institution-contingent inverted-U
%       in the internationalization-performance relationship
%       (full Asia-Pacific sample, 49 economies, WBES)
% WORKING DRAFT submission framework (Springer-compatible).
%
% COMPILE:
%   pdflatex p7_multicountry.tex   (run twice)
%   bibtex   p7_multicountry        (bibliography is embedded via filecontents)
%   pdflatex p7_multicountry.tex
%   pdflatex p7_multicountry.tex
% NOTE: No pdflatex was available in the build environment, so this file
%   ships uncompiled. It uses the standard `article` class (NOT Springer's
%   sn-jnl class, which is not bundled here) per the note in
%   templates/springer_paper.tex. To submit to a Springer Nature journal,
%   swap \documentclass and the title block for sn-jnl.cls when available.
%
% INTEGRITY NOTE: Every numeric value below is taken verbatim from the
%   dissertation (chuong_4_ket_qua_vi.md sec 4.6 / 4.9 and
%   phu_luc_A_hop_nhat_du_lieu_vi.md). Full regression coefficient tables
%   that do NOT yet exist in the dissertation are left as explicit
%   PLACEHOLDERS marked "[Table X -- full coefficients from the authors'
%   original output]". No new coefficients have been invented.
% =====================================================================

\documentclass[12pt,a4paper]{article}

\usepackage[margin=2.5cm]{geometry}
\usepackage{setspace}\onehalfspacing
\usepackage{iftex}
\ifPDFTeX
  \usepackage[utf8]{inputenc}
  \usepackage[T1]{fontenc}
  \usepackage{lmodern}
\else
  \usepackage{fontspec}
\fi
\usepackage{amsmath,amssymb,amsfonts,mathtools}
\usepackage{booktabs}
\usepackage{array}
\usepackage{tabularx}
\usepackage{graphicx}
\usepackage{float}
\usepackage{caption}\captionsetup{font=small,labelfont=bf}
\usepackage{microtype}
\usepackage{textcomp}
\usepackage[hidelinks,breaklinks=true]{hyperref}
\usepackage[round]{natbib}
\bibliographystyle{apalike}

% Embedded APA-style bibliography (subset relevant to P7).
\usepackage{filecontents}
\begin{filecontents*}{p7refs.bib}
@article{deaton1985,
  author={Deaton, Angus}, title={Panel data from a time series of cross-sections},
  journal={Journal of Econometrics}, volume={30}, number={1--2}, pages={109--126}, year={1985}}
@incollection{verbeek2008,
  author={Verbeek, Marno}, title={Pseudo-panels and repeated cross-sections},
  booktitle={The Econometrics of Panel Data}, edition={3rd}, pages={369--383},
  publisher={Springer}, year={2008}}
@article{lindmehlum2010,
  author={Lind, Jo Thori and Mehlum, Halvor},
  title={With or without {U}? The appropriate test for a {U}-shaped relationship},
  journal={Oxford Bulletin of Economics and Statistics}, volume={72}, number={1},
  pages={109--118}, year={2010}}
@article{lubeamish2004,
  author={Lu, Jane W. and Beamish, Paul W.},
  title={International diversification and firm performance: The {S}-curve hypothesis},
  journal={Academy of Management Journal}, volume={47}, number={4}, pages={598--609}, year={2004}}
@article{contractor2012,
  author={Contractor, Farok J.},
  title={Why do multinational firms exist? A theory note about the effect of multinational
  expansion on performance and recent methodological critiques},
  journal={Global Strategy Journal}, volume={2}, number={4}, pages={318--331}, year={2012}}
@article{marano2016,
  author={Marano, Valentina and Arregle, Jean-Luc and Hitt, Michael A. and Spadafora, Ettore
  and Van Essen, Marc},
  title={Home country institutions and the internationalization-performance relationship:
  A meta-analytic review},
  journal={Journal of Management}, volume={42}, number={5}, pages={1075--1110}, year={2016}}
@article{kirca2012,
  author={Kirca, Ahmet H. and Roth, Kendall and Hult, G. Tomas M. and Cavusgil, S. Tamer},
  title={The role of context in the multinationality-performance relationship: A meta-analytic review},
  journal={Global Strategy Journal}, volume={2}, number={2}, pages={108--121}, year={2012}}
@article{kaufmann2011,
  author={Kaufmann, Daniel and Kraay, Aart and Mastruzzi, Massimo},
  title={The Worldwide Governance Indicators: Methodology and analytical issues},
  journal={Hague Journal on the Rule of Law}, volume={3}, number={2}, pages={220--246}, year={2011}}
@book{khanna2010,
  author={Khanna, Tarun and Palepu, Krishna G.},
  title={Winning in Emerging Markets: A Road Map for Strategy and Execution},
  publisher={Harvard Business Press}, year={2010}}
@book{north1990,
  author={North, Douglass C.}, title={Institutions, Institutional Change and Economic Performance},
  publisher={Cambridge University Press}, year={1990}}
@article{solon2015,
  author={Solon, Gary and Haider, Steven J. and Wooldridge, Jeffrey M.},
  title={What are we weighting for?},
  journal={Journal of Human Resources}, volume={50}, number={2}, pages={301--316}, year={2015}}
@article{cameron2015,
  author={Cameron, A. Colin and Miller, Douglas L.},
  title={A practitioner's guide to cluster-robust inference},
  journal={Journal of Human Resources}, volume={50}, number={2}, pages={317--372}, year={2015}}
@article{aterido2011,
  author={Aterido, Reyes and Hallward-Driemeier, Mary and Pag\'es, Carmen},
  title={Big constraints to small firms' growth? Business environment and employment growth across firms},
  journal={Economic Development and Cultural Change}, volume={59}, number={3}, pages={609--647}, year={2011}}
@article{richard2009,
  author={Richard, Pierre J. and Devinney, Timothy M. and Yip, George S. and Johnson, Gerry},
  title={Measuring organizational performance: Towards methodological best practice},
  journal={Journal of Management}, volume={35}, number={3}, pages={718--804}, year={2009}}
\end{filecontents*}

\title{Institution-Contingent Turning Points in the
Internationalization--Performance Relationship: Evidence from 49
Asia--Pacific Economies}
\author{Author details withheld for blind review}
\date{Working draft}

\begin{document}
\maketitle

\begin{abstract}
\noindent The internationalization--performance (I--P) relationship is
widely theorized to follow an inverted-U, yet meta-analyses show that the
location of its turning point is unstable across contexts. We test whether
the formal-institution quality of the home economy systematically shifts
the I--P turning point, using a pooled repeated cross-section of World Bank
Enterprise Survey establishments spanning 49 Asia--Pacific economies. Labour
productivity is standardized within country--year to neutralize currency
artefacts, and all models carry two-way (country and year) fixed effects so
the curvature is identified from within-country variation. In the strongest
specification ($N=38{,}342$; adjusted $R^2 = 0.677$) the pooled turning
point is $40.0\%$ of foreign sales to total sales, robust across all
specifications (turning point range $33.8$--$40.0\%$) and reconfirmed in the
full three-way moderation model ($34.6\%$). The central result is an
institution-contingent gradient: the turning point rises monotonically from
roughly $28\%$ in advanced-innovation economies to roughly $55\%$ in
emerging and small-island economies, mirroring an independent meta-analytic
moderation test ($Q_M = 17.35$, $p = 0.002$). Digital adoption flattens and
lifts the curve, with its protective role strongest where institutions are
weakest. Findings reconcile the conflicting I--P literature by reframing the
turning point as an institutionally-conditioned, rather than universal,
parameter.\\[4pt]
\textbf{Keywords:} internationalization; firm performance; inverted-U;
institutional quality; turning point; emerging markets; Asia--Pacific
\end{abstract}

\section{Introduction}
The relationship between internationalization and firm performance is one of
the most studied yet least settled questions in international business.
Decades of work document inverted-U, S-shaped, and linear forms, and
meta-analyses attribute much of this disagreement to unobserved contextual
moderators rather than to genuine theoretical disagreement
\citep{kirca2012,marano2016,contractor2012}. A recurring gap is that most
primary studies estimate a single curve for a single economy, leaving the
\emph{location} of the turning point---the degree of internationalization at
which marginal returns turn negative---uncompared across institutional
settings.

This study addresses that gap directly. We assemble a large, harmonized,
multi-country establishment sample and ask whether home-economy institutional
quality systematically relocates the I--P turning point. Our contribution is
threefold. First, we provide what is, to our knowledge, the largest
single-design primary test of the inverted-U across Asia--Pacific economies,
estimated within a unified harmonization protocol. Second, we document an
\emph{institution-contingent gradient}: the turning point moves outward as
institutional quality declines, a pattern predicted by institutional theory
but rarely shown with primary firm-level data. Third, the primary gradient
converges with an independent meta-analytic moderation result, giving the
finding cross-method robustness.

\section{Theory and Hypotheses}
The inverted-U rests on a benefits-then-costs logic: early international
expansion yields learning and scale economies, while further expansion incurs
rising coordination, liability-of-foreignness, and governance costs
\citep{lubeamish2004,contractor2012}. Institutional theory implies the
\emph{balance point} of this trade-off is conditioned by the home
institutional environment \citep{north1990,khanna2010,marano2016}. Where
formal institutions are strong, firms recover the costs of expansion at lower
exposure and so reach peak performance earlier (an earlier, lower turning
point). Where institutions are weak, firms must internationalize more heavily
to offset domestic transaction costs and institutional voids, pushing the
turning point outward.

\begin{description}
\item[H1.] The I--P relationship is inverted-U shaped across the pooled
Asia--Pacific sample.
\item[H5.] The turning point of the inverted-U increases monotonically as
home-economy institutional quality declines (the ICRV gradient).
\item[H3.] Digital adoption (DAI) moderates the curvature of the I--P
relationship, with stronger protective effects in weaker institutional
contexts.
\end{description}

\section{Data and Methods}
\subsection{Data source and unit of analysis}
The primary data are World Bank Enterprise Surveys (WBES), establishment-level
surveys collected under a standardized global instrument with stratified
random sampling by sector, size, and region \citep{aterido2011}. The unit of
analysis is the establishment. Because WBES draws an \emph{independent} sample
each wave, the data are repeated cross-sections rather than a balanced panel
\citep{deaton1985,verbeek2008}; only waves from 2006 onward (the unified
global methodology) are retained.

\subsection{Harmonization and analytic sample}
A six-step pooling and harmonization protocol (S1--S6; full detail in the
data appendix) screens survey types, de-duplicates country--year pairs, maps
heterogeneous fields to a common schema, recodes WBES non-response codes
($-9$/$-8$/$-7$) to missing, assigns institutional regimes, and applies
listwise deletion on the focal variables. The dataflow yields a classified
pool of $96{,}415$ firms (52 economies), an Asia--Pacific analytic frame of
$91{,}982$ firms across 49 economies, a baseline quadratic regression sample
of $N = 84{,}910$ (model M2), and a fully-controlled sample of $N = 38{,}342$
(models M3--M5). [These sample sizes are REAL, taken from Appendix~A,
Step~S5--S6.]

\subsection{Measures}
Labour productivity (the dependent variable) is $\ln(\text{annual
sales}/\text{full-time employees})$, standardized ($z$) within each
country--year to remove currency and price-level artefacts
\citep{richard2009}. Internationalization intensity (FSTS) is the share of
direct plus indirect exports in sales, bounded $[0,100]$. Technological
capability (TCI) and digital adoption (DAI) are built from binary indicators
(international quality certification, foreign-licensed technology, website,
e-mail, electronic payment).

\subsection{Estimation}
We estimate a pooled cross-section with two-way fixed effects:
\begin{equation}
\ln(\mathrm{LP})_{ijt} = \beta_1 \mathrm{FSTS}_{ijt}
  + \beta_2 \mathrm{FSTS}^2_{ijt}
  + \mathbf{X}_{ijt}\boldsymbol{\gamma}
  + \alpha_j + \tau_t + \varepsilon_{ijt},
\end{equation}
where $\alpha_j$ and $\tau_t$ are economy and year fixed effects. Economy
fixed effects absorb all country-invariant heterogeneity (currency, price
level, baseline institutions), so the curvature is identified from
within-country variation \citep{solon2015}. The turning point is
$\mathrm{TP} = -\beta_1/(2\beta_2)$, tested with the Lind--Mehlum (LM)
procedure. Standard errors are heteroskedasticity-robust and clustered by
economy \citep{cameron2015}.

\section{Results}
\subsection{Pooled inverted-U and turning point (H1)}
The strongest specification, M5 (country and year fixed effects;
adjusted $R^2 = 0.677$), yields a pooled turning point of
\[
\mathrm{TP(M5)} = 40.0\%\ \mathrm{FSTS}\qquad (p_{\mathrm{LM}} < 0.001).
\]
The turning point stays within a narrow band of $33.8$--$40.0\%$ FSTS across
all specifications M2--M9, and is reconfirmed in the full three-way
moderation model M11 ($\mathrm{TP}=34.6\%$, $p_{\mathrm{LM}}=0.002$,
$N=29{,}840$). This consistency rules out the turning point being an artefact
of small samples or omitted controls. [All values in this paragraph are REAL.]

% --------------------------------------------------------------------
% PLACEHOLDER: full coefficient table not present in the dissertation.
\begin{table}[H]
\centering
\caption{Quadratic FSTS specifications M2--M5 (full coefficients).}
\begin{tabularx}{\textwidth}{Xcccc}
\toprule
 & M2 & M3 & M4 & M5 \\
\midrule
\multicolumn{5}{c}{\itshape [Table P7-1 -- full coefficients from the authors'
original output (do-file). Do NOT populate without the locked estimates.]}\\
FSTS              & --- & --- & --- & --- \\
FSTS$^2$          & --- & --- & --- & --- \\
Turning point     & --- & --- & --- & $40.0\%$ \\
$N$               & $84{,}910$ & --- & --- & $38{,}342$ \\
Adjusted $R^2$    & --- & --- & --- & $0.677$ \\
Country FE / Year FE & --- & --- & --- & Yes / Yes \\
\bottomrule
\end{tabularx}
\end{table}
% --------------------------------------------------------------------

\subsection{The institution-contingent gradient (H5)}
The central result is a monotone ICRV gradient in the turning point:
\[
\mathrm{TP}(\text{Advanced Innovation}) \approx 28\%
  < \mathrm{TP}(\text{Upper-middle})
  < \mathrm{TP}(\text{Emerging/SIDS}) \approx 55\%.
\]
Firms in strong-institution economies reach peak productivity at lower
internationalization; firms in weak-institution economies require heavier
internationalization to offset transaction costs and institutional
uncertainty, so the turning point appears later. An overall ICRV moderation
test is highly significant and converges with the dissertation's
meta-analysis ($Q_M = 17.35$, $p = 0.002$), giving independent confirmation
across methods. [Gradient bounds $28\%$/$55\%$ and $Q_M$ are REAL.]

\subsection{Capability and digital moderation (H3)}
TCI raises the productivity level across the sample (a one-SD increase lifts
labour productivity by about $41\%$, $\exp(0.344)-1$), but its
curvature interactions are not significant in the 49-economy sample, so the
capability-shaping role of TCI is context-specific rather than universal. DAI
shows a significant level effect ($\hat{\beta}(\mathrm{DAI}) = +0.155$,
$p < 0.001$; about a $17\%$ productivity advantage for firms with a website)
and significant curvature interactions
($\hat{\beta}(\mathrm{FSTS}\times\mathrm{DAI}) = -0.614$, $p<0.001$;
$\hat{\beta}(\mathrm{FSTS}^2\times\mathrm{DAI}) = +0.766$, $p<0.01$). The
DAI$\times$ICRV interaction ($p = 0.012$) indicates the digital ``shield'' is
strongest where institutions are weakest. [All coefficients in this
paragraph are REAL.]

\begin{figure}[H]
\centering
% Figure ships with the dissertation; path is relative to thesis/figures.
% \includegraphics[width=.8\textwidth]{fig_4_x_ip_curves_icrv_gradient.png}
\fbox{\parbox{.8\textwidth}{\centering\vspace{2.2cm}
[Figure~1 -- family of inverted-U I--P curves with turning point shifting
rightward as institutional quality declines (28\% to 55\%); SIDS curve turns
monotone negative. Source figure:
thesis/figures/fig\_4\_x\_ip\_curves\_icrv\_gradient.png]\vspace{2.2cm}}}
\caption{Institution-contingent family of I--P curves (ICRV gradient).}
\end{figure}

\section{Discussion and Conclusion}
The evidence reframes the long-standing I--P controversy. The inverted-U is
robust, but its turning point is not a universal constant; it is an
institutionally-conditioned parameter that moves outward as home-economy
institutional quality falls. This reconciles inconsistent prior findings:
studies set in different institutional contexts were, in effect, sampling
different points of a single contingent family of curves. The convergence of
the primary gradient with an independent meta-analytic moderation test
strengthens the inference. Digital adoption emerges as a partial substitute
for institutional support, flattening and lifting the curve most where
institutions are weakest.

Limitations follow from the design. Repeated cross-sections preclude firm
trajectory tracking and dynamic causal inference; labour productivity is a
proxy for performance; and listwise deletion reduces $N$ as controls are
added. A data update adding Japan-2025 and new 2024--2026 waves is reserved
for a subsequent re-estimation round; the results reported here are held on
the locked 49-economy analytic sample. Future work should re-estimate the
gradient on the expanded frame and probe the digital-substitution mechanism
with finer measures.

\bibliography{p7refs}

\end{document}
